\documentclass[12pt]{article}

\usepackage[T1]{fontenc}
\usepackage{amsmath,amssymb,amsthm}
\usepackage{tikz-cd}
\usepackage{tikz}
\usepackage{hyperref}
\usepackage{cleveref}
\usepackage{graphicx}
\usepackage[margin=1in]{geometry}
\usepackage{everypage}
\usepackage{xcolor}
\usepackage{booktabs}
\usepackage{enumitem}
\usepackage{mathtools}
\usepackage{microtype}
\usepackage{array}
\usepackage{longtable}

\hypersetup{
  colorlinks=true,
  linkcolor=blue!55!black,
  citecolor=green!45!black,
  urlcolor=blue!65!black
}

\newtheorem{theorem}{Theorem}[section]
\newtheorem{proposition}[theorem]{Proposition}
\newtheorem{lemma}[theorem]{Lemma}
\newtheorem{corollary}[theorem]{Corollary}
\theoremstyle{definition}
\newtheorem{definition}[theorem]{Definition}
\newtheorem{example}[theorem]{Example}
\newtheorem{assumption}[theorem]{Assumption}
\theoremstyle{remark}
\newtheorem{remark}[theorem]{Remark}
\newtheorem{warning}[theorem]{Warning}

\newcommand{\A}{\mathcal A}
\newcommand{\B}{\mathcal B}
\newcommand{\K}{\mathcal K}
\newcommand{\Ham}{\mathfrak{Ham}}
\newcommand{\EFT}{\mathfrak{EFT}}
\newcommand{\Class}{\mathfrak{Class}}
\newcommand{\Real}{\mathfrak{Real}}
\newcommand{\spec}{\operatorname{spec}}
\newcommand{\sgn}{\operatorname{sgn}}
\newcommand{\dist}{\operatorname{dist}}
\newcommand{\Cont}{\operatorname{Cont}}
\newcommand{\diag}{\operatorname{diag}}
\newcommand{\id}{\operatorname{id}}
\newcommand{\im}{\operatorname{im}}
\newcommand{\status}[1]{\textsc{#1}}
\newcommand{\Z}{\mathbb Z}
\newcommand{\R}{\mathbb R}
\newcommand{\C}{\mathbb C}
\newcommand{\T}{\mathbb T}

\definecolor{grokgray}{RGB}{110,110,110}
\AddEverypageHook{%
  \ifnum\value{page}=1
    \begin{tikzpicture}[remember picture,overlay]
      \node[
        rotate=90,
        anchor=south,
        font=\Large\sffamily\bfseries\color{grokgray},
        inner sep=0pt
      ] at ([xshift=38pt,yshift=0.52\paperheight]current page.south west)
      {GrokRxiv:2026.07.microscopic-effective-theories\quad
       [\,math-ph\,]\quad14 Jul 2026};
    \end{tikzpicture}
  \fi
}

\title{From Microscopic Lattice Hamiltonians to Effective Theories:\\
Controlled K-Theory Maps, Information Loss, and Realization Limits}

\author{Matthew Long\\
\textit{The YonedaAI Collaboration}\\
\textit{YonedaAI Research Collective}\\
Chicago, IL\\
\texttt{matthew@yonedaai.com} $\cdot$ \url{https://yonedaai.com}}

\date{14 July 2026}

\begin{document}

\maketitle

\begin{abstract}
We separate three maps that are often compressed into a claim of
``microscopic and effective equivalence.''  The first sends a specified
microscopic Hamiltonian to a low-energy description, when a scaling or
spectral reduction exists.  The second sends controlled low-energy data to a
stable homotopy, operator K-theory, or bordism class.  The third asks for an
explicit local lattice representative of an abstract class.  No theorem
identifies all three maps for general interacting matter.

Our strongest result concerns bounded free-fermion Hamiltonians in a unital
C*-algebra.  A common Fermi-level gap permits continuous functional calculus,
spectral flattening, a Fermi projection, and a relative operator K-class.
We prove naturality in compact parameter spaces, homotopy invariance, and
additivity under stacking.  We then give an information-loss contract: the
class retains stable occupied-sector and symmetry data, while it discards gap
magnitude, dispersion, velocity, locality constants, and microscopic
couplings.  Covariant disordered systems fit the same construction through a
crossed-product algebra when a genuine spectral gap is present.

The Su-Schrieffer-Heeger family supplies a complete finite example.  Its
gap, flattening, winding difference, and local transition charge are computed
explicitly.  A staggered sublattice mass gives a gapped detour between the two
chiral phases once chiral symmetry is dropped.  The Fidkowski-Kitaev reduction
of the one-dimensional BDI free classification from $\Z$ to $\Z/8$ under
interactions supplies a different boundary: free operator K-theory cannot be
promoted unchanged to a universal interacting classification.  Condensed
parameter objects organize continuous and profinite families, but they do not
create a scaling limit, prove an EFT equivalence, or solve the inverse
realization problem.
\end{abstract}

\tableofcontents

\section{Introduction}

The phrase ``the effective theory of a lattice model'' can refer to several
different constructions.
A Bloch Hamiltonian may be linearized near a band crossing.
A gapped one-particle Hamiltonian may be flattened to an involution.
A many-body system may have a continuum scaling limit.
An invertible low-energy theory may determine a bordism class.
These operations have different hypotheses and forget different data.

This paper treats the fourth problem in a program on topological phases in a
condensed-mathematics setting:

\[
\begin{tikzcd}[row sep=large]
\text{local quantum interactions}
  \arrow[d]
\\
\text{condensed moduli stack of Hamiltonians}
  \arrow[d]
\\
\text{uniformly gapped substack}
  \arrow[d]
\\
\text{stabilized phase }\infty\text{-groupoid}
  \arrow[d]
\\
\text{invertible condensed phase spectrum.}
\end{tikzcd}
\]

The earlier analytic steps remain input here.
Locality supplies a quasi-local observable algebra and controlled dynamics.
Positivity and the C*-norm make functional calculus meaningful.
A thermodynamic gap supplies the physical phase condition.
The present question is what can be transported from a microscopic system to
an effective or homotopical target, and whether that transport can be
reversed.

\subsection{The three arrows}

We insist on the following factorization:

\[
\begin{tikzcd}[column sep=large]
\Ham_{\mathrm{micro}}
  \arrow[r,"L"]
&
\EFT_{\mathrm{low}}
  \arrow[r,"I"]
&
\Class_{\mathrm{stable}}
\end{tikzcd}
\qquad
\begin{tikzcd}
\Class_{\mathrm{stable}}
  \arrow[r,dashed,"R"]
&
\Real_{\mathrm{lattice}}.
\end{tikzcd}
\]

Here $L$ is a low-energy or spectral-reduction map.
The map $I$ extracts a stable invariant.
The dashed map $R$ is a realization problem.
It is not a formal inverse of $I\circ L$.

\begin{definition}[Microscopic-to-low-energy comparison]
A microscopic-to-low-energy comparison consists of a domain
$\mathcal D\subset\Ham_{\mathrm{micro}}$, a target category
$\EFT_{\mathrm{low}}$, and a rule
\[
L:\mathcal D\longrightarrow\EFT_{\mathrm{low}}
\]
together with a stated limit, approximation, or functional-calculus theorem
that controls the rule.
\end{definition}

\begin{definition}[Invariant extraction]
An invariant extraction is a map
\[
I:\EFT_{\mathrm{low}}\longrightarrow\Class_{\mathrm{stable}}
\]
defined after fixing dimension, symmetry, stabilization, and equivalence.
It is an equivalence only if faithfulness, fullness, and essential
surjectivity are separately proved.
\end{definition}

\begin{definition}[Microscopic realization]
A realization of $c\in\Class_{\mathrm{stable}}$ is an explicit local or
uniformly almost-local interaction, symmetry action, ground-state convention,
and system-size-independent gap witness whose comparison class is $c$.
\end{definition}

The definitions prevent a common reversal error.
The existence of an invariant map does not provide a microscopic model for
every target class.
Agreement of invariants does not imply a microscopic gapped path unless a
completeness theorem has been proved in the chosen domain.

\subsection{Contributions}

The paper makes six claims of different strength.

\begin{enumerate}[label=\arabic*.]
  \item We prove a controlled functional-calculus theorem for compact
  families of bounded free-fermion Hamiltonians with a common Fermi gap.

  \item We define a relative operator K-class with respect to a chosen
  reference Hamiltonian and prove homotopy invariance and stacking
  additivity.

  \item We give a precise retention and loss table for spectral flattening.
  Positive energy rescaling supplies a simple nonfaithfulness proof.

  \item We extend the controlled construction to covariant disordered
  one-particle systems represented in
  $C(\Omega)\rtimes\Z^d$, under a genuine C*-spectral gap.

  \item We compute the SSH transition charge on a linking circle and prove
  that a staggered mass provides a gapped symmetry-breaking detour.

  \item We use the BDI interaction reduction to show why the free result has
  a sharp domain boundary.
\end{enumerate}

The paper does not claim a general microscopic/EFT equivalence.
It also does not claim that every bordism or homotopy class has a local
gapped lattice representative.

\subsection{Status vocabulary}

Every substantive result is assigned one of the following labels.

\begin{center}
\begin{tabular}{p{0.18\textwidth}p{0.70\textwidth}}
\toprule
Label & Meaning \\
\midrule
\status{established} & A cited theorem or standard construction under the assumptions displayed here.\\
\status{proposed} & A definition or condensed organization introduced by this program.\\
\status{conjectural} & A precise assertion for which this paper gives no proof.\\
\status{open} & No proof or counterexample is known in the stated scope.\\
\status{obstructed} & A theorem or explicit counterexample rules out the claim as stated.\\
\bottomrule
\end{tabular}
\end{center}

The functional-calculus statements below are \status{established}.
The condensed packaging is \status{proposed}.
A general lattice/EFT equivalence is \status{open}.
An unchanged extension of free K-theory to all interacting BDI phases is
\status{obstructed}.

\section{Domains, targets, and equivalences}

\subsection{Microscopic quadratic systems}

Let $\Gamma$ be a countable metric lattice and let $\mathcal H$ be the
one-particle Hilbert space.
In the simplest translation-invariant case,
\[
\mathcal H=\ell^2(\Z^d)\otimes\C^N.
\]
A finite-range one-particle Hamiltonian is a bounded self-adjoint operator
whose matrix elements vanish beyond a fixed lattice distance.
Exponential or summable decay gives broader controlled classes.

After Fourier transform, a translation-invariant finite-range Hamiltonian is
a matrix-valued trigonometric polynomial
\[
h:\T^d\longrightarrow M_N(\C),
\qquad
h(k)^*=h(k).
\]
More generally, one works with a self-adjoint element
\[
h\in M_N(\A)
\]
of a unital C*-algebra $\A$ encoding position, covariance, disorder, or
symmetry.

\begin{definition}[Controlled free-fermion datum]
A controlled free-fermion datum is a tuple
\[
(\A,N,h,\mu,\mathcal S,h_{\mathrm{ref}})
\]
where $\A$ is a specified real, complex, graded, or twisted C*-algebra,
$h=h^*\in M_N(\A)$, $\mu\in\R$ is a Fermi level,
$\mathcal S$ records the symmetry constraints, and
$h_{\mathrm{ref}}$ is a reference datum of the same type.
The datum is spectrally gapped when
\[
\dist(\mu,\spec(h))>0.
\]
\end{definition}

We set $\mu=0$ after replacing $h$ by $h-\mu$.
The algebra and the representation class are part of the problem.
Changing them can change the relevant K-group.

\subsection{Interacting lattice systems}

An interacting microscopic system is not generally described by a bounded
element $h\in\A_\Gamma$.
It is described by an interaction
\[
\Phi:X\Subset\Gamma\longmapsto\Phi(X)\in\A_X^{\mathrm{sa}},
\]
finite-volume Hamiltonians, and thermodynamic dynamics.
The formal infinite sum of all $\Phi(X)$ need not converge in the quasi-local
algebra.

This distinction matters for spectral flattening.
Continuous functional calculus applies directly to a bounded self-adjoint
element of a C*-algebra.
It does not turn an extensive many-body Hamiltonian into a bounded
quasi-local observable.
Many-body phase equivalence instead uses a uniform finite-volume or bulk gap,
quasi-adiabatic continuation, and stabilization by atomic systems.

\begin{warning}[Domain boundary]
The theorem in \cref{sec:functional-calculus} concerns bounded one-particle
Hamiltonians or bounded C*-algebra representatives.
It is not a spectral-flattening theorem for an arbitrary interacting
infinite-volume Hamiltonian.
\end{warning}

\subsection{What counts as an effective theory}

The term effective field theory will be used in a restricted way.
An effective description records fields, symmetries, couplings, and a regime
of momenta or energies in which its correlation functions approximate those
of a microscopic model to a stated accuracy.
A topological field theory is a still coarser target, usually intended to
retain only long-distance invertible response or defect data.

\begin{definition}[Controlled effective description]
For a microscopic datum $H$, a controlled effective description consists of
a scale $\Lambda_{\mathrm{IR}}$, an effective action or Hamiltonian
$S_{\mathrm{eff}}$, observables $\mathcal O$, and an error statement
\[
\left|
\langle\mathcal O\rangle_H
-
\langle\mathcal O\rangle_{S_{\mathrm{eff}}}
\right|
\leq \varepsilon(\Lambda_{\mathrm{IR}})
\]
for a specified class of observables, with
$\varepsilon(\Lambda_{\mathrm{IR}})\to0$ in a specified limit.
\end{definition}

Spectral flattening is not such an approximation.
It is a homotopy inside the space of gapped bounded operators.
It preserves a stable topological class while deliberately destroying the
energy scale.

\subsection{An equivalence audit}

Any claimed comparison must answer more than whether a map can be written.
We use the following audit throughout.

\begin{longtable}{p{0.22\textwidth}p{0.31\textwidth}p{0.35\textwidth}}
\toprule
Question & Required datum & Failure mode \\
\midrule
\endfirsthead
\toprule
Question & Required datum & Failure mode \\
\midrule
\endhead
Construction & A defined map on a stated domain & Formal analogy with no map \\
Homotopy invariance & A chosen path topology and gap condition & Gap closes along the path \\
Faithfulness & Distinct inputs remain distinct & Flattening forgets energy scales \\
Fullness & Target morphisms lift to microscopic morphisms & EFT deformation has no lattice lift \\
Essential surjectivity & Every target object has a preimage & Abstract class lacks a local model \\
Stacking & Direct sum or tensor product is respected & Stabilization conventions disagree \\
Symmetry & Real, graded, or twisted structure is retained & Forgetting symmetry opens a detour \\
Boundaries and defects & Boundary conditions enter the category & Bulk class omits boundary choices \\
Interactions & Domain is closed under allowed interactions & Free classification collapses \\
Uniform gap & One lower bound controls the family & Pointwise gaps approach zero \\
\bottomrule
\end{longtable}

Calling a comparison an equivalence requires affirmative answers in the
chosen category.
The controlled free-fermion map below proves construction, gapped-homotopy
invariance, and stacking compatibility.
It is intentionally not faithful to microscopic spectra.

\section{Functional calculus and the controlled K-class}
\label{sec:functional-calculus}

\subsection{A common spectral gap}

Let $S$ be compact Hausdorff and $\A$ a unital complex C*-algebra.
Write
\[
\A_S=C(S,\A).
\]
An element $h\in M_N(\A_S)$ is equivalently a norm-continuous family
$s\mapsto h_s\in M_N(\A)$.

\begin{definition}[Uniform Fermi gap]
A self-adjoint family $h=h^*\in M_N(\A_S)$ has a uniform Fermi gap
$\gamma>0$ when
\[
\spec(h_s)\cap(-\gamma,\gamma)=\varnothing
\]
for every $s\in S$.
\end{definition}

Because $S$ is compact and $h$ is norm-continuous, pointwise invertibility is
enough to produce some common positive gap.
We state the quantitative constant because later parameter objects need not
be compact, and because estimates depend on it.

\begin{lemma}[Inverse and gap]
\label{lem:inverse-gap}
For self-adjoint $a$ in a unital C*-algebra, the following are equivalent:
\begin{enumerate}[label=(\roman*)]
  \item $0\notin\spec(a)$;
  \item $a$ is invertible;
  \item there is $\gamma>0$ with
  $\spec(a)\cap(-\gamma,\gamma)=\varnothing$.
\end{enumerate}
For an invertible self-adjoint $a$,
\[
\dist(0,\spec(a))=\|a^{-1}\|^{-1}.
\]
\end{lemma}

\begin{proof}
The equivalence of the first two statements is the definition of the
spectrum.
The spectrum is compact, so exclusion of zero is equivalent to a positive
distance from zero.
For a normal element, continuous functional calculus gives
\[
\|a^{-1}\|
=
\sup_{\lambda\in\spec(a)}|\lambda|^{-1},
\]
which proves the last identity.
\end{proof}

\subsection{Spectral flattening}

Define the sign function on $\R\setminus\{0\}$ by
\[
\sgn(x)=
\begin{cases}
1,&x>0,\\
-1,&x<0.
\end{cases}
\]
It is continuous on the spectrum of any invertible self-adjoint element.

\begin{theorem}[Controlled spectral flattening]
\label{thm:flattening}
Let $S$ be compact Hausdorff, let $\A$ be a unital complex C*-algebra, and
let $h=h^*\in M_N(C(S,\A))$ have a uniform Fermi gap.
Set
\[
q=\sgn(h)=h|h|^{-1},
\qquad
p=\frac{1-q}{2}.
\]
Then:
\begin{enumerate}[label=(\alph*)]
  \item $q=q^*$ and $q^2=1$;
  \item $p=p^*=p^2$;
  \item evaluation commutes with flattening,
  $q(s)=\sgn(h_s)$ and $p(s)=\chi_{(-\infty,0)}(h_s)$;
  \item the path
  \[
  h_u=(1-u)h+u\,q,
  \qquad 0\leq u\leq1,
  \]
  remains invertible and preserves the positive and negative spectral
  subspaces.
\end{enumerate}
\end{theorem}

\begin{proof}
Continuous functional calculus in $M_N(C(S,\A))$ defines
$q=\sgn(h)$.
The scalar identities
$\sgn(x)^2=1$ and $\sgn(x)=\overline{\sgn(x)}$ on $\spec(h)$ imply
$q^2=1$ and $q=q^*$.
The projection identities for $p$ follow by direct algebra.

Every evaluation map
\[
\operatorname{ev}_s:C(S,\A)\longrightarrow\A
\]
is a unital star-homomorphism.
Functoriality of continuous functional calculus gives part (c).

For $\lambda\in\spec(h)$, the scalar homotopy is
\[
f_u(\lambda)=(1-u)\lambda+u\,\sgn(\lambda).
\]
If $\lambda>0$, then $f_u(\lambda)>0$.
If $\lambda<0$, then $f_u(\lambda)<0$.
Thus zero never enters the spectral image.
Functional calculus identifies $h_u=f_u(h)$ and proves part (d).
\end{proof}

\begin{remark}
The theorem is sometimes described as replacing all occupied energies by
$-1$ and all empty energies by $+1$.
That description is correct only after the Fermi level and occupied
convention are fixed.
\end{remark}

\subsection{Relative operator K-theory}

The projection $p$ determines a class
\[
[p]\in K_0(C(S,\A)).
\]
Following the difference-group viewpoint emphasized by Thiang
\cite{Thiang2016}, we compare it with a reference system.

\begin{definition}[Relative difference class]
Let $h$ and $h_{\mathrm{ref}}$ be uniformly gapped self-adjoint elements of
matrix algebras over $C(S,\A)$, with all required symmetry data fixed.
Their relative class is
\[
\kappa_S(h,h_{\mathrm{ref}})
=
[p_h]-[p_{h_{\mathrm{ref}}}]
\in K_0(C(S,\A)).
\]
\end{definition}

The word relative here refers to a difference from a reference phase.
It should not be confused with the relative group $K^*(B,U)$ of a topological
pair, which appears in the transition calculation later.

\begin{proposition}[Gapped-homotopy invariance]
\label{prop:homotopy-invariance}
Suppose $t\mapsto h_t$ is a norm-continuous path of self-adjoint elements in
$M_N(C(S,\A))$ and one $\gamma>0$ satisfies
\[
\spec(h_t(s))\cap(-\gamma,\gamma)=\varnothing
\]
for every $(t,s)\in[0,1]\times S$.
Then $t\mapsto p_{h_t}$ is a norm-continuous path of projections and
\[
[p_{h_0}]=[p_{h_1}]
\quad\text{in}\quad K_0(C(S,\A)).
\]
Consequently $\kappa_S(h_t,h_{\mathrm{ref}})$ is constant in $t$.
\end{proposition}

\begin{proof}
Regard the family as one self-adjoint element
\[
H\in M_N(C([0,1]\times S,\A)).
\]
The common gap makes $H$ invertible.
Theorem \ref{thm:flattening} gives a projection
$P=\chi_{(-\infty,0)}(H)$.
Evaluation at $t$ yields $p_{h_t}$.
Thus the endpoints are homotopic projections and define the same K-class.
\end{proof}

\begin{proposition}[Stacking additivity]
\label{prop:stacking}
For two controlled data $h_1,h_2$ and references
$r_1,r_2$ over the same parameter space,
\[
\kappa_S(h_1\oplus h_2,r_1\oplus r_2)
=
\kappa_S(h_1,r_1)+\kappa_S(h_2,r_2).
\]
\end{proposition}

\begin{proof}
Functional calculus respects block direct sums, so
\[
p_{h_1\oplus h_2}=p_{h_1}\oplus p_{h_2}.
\]
Addition in $K_0$ is defined by block direct sum.
Subtracting the corresponding reference identity gives the formula.
\end{proof}

\subsection{Chiral symmetry and the odd class}

Suppose a grading
\[
\Gamma=
\begin{pmatrix}
1&0\\
0&-1
\end{pmatrix}
\]
satisfies $\Gamma h\Gamma=-h$.
Functional calculus preserves this oddness:
$\Gamma q\Gamma=-q$.

\begin{proposition}[Off-diagonal unitary]
\label{prop:chiral-unitary}
Assume the two graded summands have equal stabilized size.
Then the flattened operator has the form
\[
q=
\begin{pmatrix}
0&u^*\\
u&0
\end{pmatrix}
\]
for a unitary $u$ over $C(S,\A)$.
The chiral stable class is $[u]\in K_1(C(S,\A))$, or the relative difference
$[u]-[u_{\mathrm{ref}}]$ after choosing a reference.
\end{proposition}

\begin{proof}
Oddness forces the diagonal blocks of $q$ to vanish.
Self-adjointness identifies the upper-right block with the adjoint of the
lower-left block.
The identity $q^2=1$ then gives
$u^*u=uu^*=1$.
The standard unitary picture of $K_1$ supplies the class.
\end{proof}

Real symmetries and particle-hole structures require the appropriate real or
graded K-group.
The complex formula must not be reused after discarding those structures.
Kitaev's periodic table \cite{Kitaev2009} and the operator-algebraic treatment
of Thiang \cite{Thiang2016} explain these symmetry-sensitive targets.

\subsection{Naturality in the parameter space}

Let $f:T\to S$ be continuous between compact Hausdorff spaces.
Pullback gives a star-homomorphism
\[
f^*:C(S,\A)\longrightarrow C(T,\A),
\qquad
(f^*a)(t)=a(f(t)).
\]

\begin{proposition}[Parameter naturality]
\label{prop:naturality}
Under the assumptions of \cref{thm:flattening},
\[
p_{f^*h}=f^*p_h,
\qquad
\kappa_T(f^*h,f^*h_{\mathrm{ref}})
=
f^*\kappa_S(h,h_{\mathrm{ref}}).
\]
\end{proposition}

\begin{proof}
The first identity is functoriality of continuous functional calculus under
the star-homomorphism $f^*$.
The induced map on $K_0$ is additive and sends differences to differences,
which proves the second identity.
\end{proof}

This naturality is the reliable part of the condensed interpretation.
For a representable condensed parameter object $\underline S$, one evaluates
the family on a compact or profinite test space and pulls back along maps of
tests.
Extending the construction to arbitrary condensed anima requires a sheaf or
stack model of the Hamiltonian domain and is \status{proposed}, not proved by
the elementary theorem above.

\section{The information-loss contract}

\subsection{What the class retains}

Within the controlled domain, spectral flattening retains the following
stable data.

\begin{enumerate}[label=(\roman*)]
  \item It retains the occupied versus empty grading determined by the chosen
  Fermi level.

  \item It retains the connected component of the gapped self-adjoint
  operator in the stabilized topology.

  \item It retains symmetry constraints that were built into the real,
  graded, or twisted target algebra.

  \item It retains the direct-sum law represented by addition in K-theory.

  \item Through pairings with cyclic cocycles or geometric cycles, it may
  retain quantized response coefficients under the additional hypotheses of
  those index theorems \cite{Bellissard1994,ProdanSchulzBaldes2016}.
\end{enumerate}

The last item is not automatic from the abstract K-class alone.
One must supply the pairing, regularity algebra, trace, and covariance data.

\subsection{What the class forgets}

The same map forgets several physically meaningful quantities.

\begin{center}
\begin{tabular}{p{0.31\textwidth}p{0.55\textwidth}}
\toprule
Discarded datum & Reason \\
\midrule
Gap magnitude & Every nonzero eigenvalue is moved to $+1$ or $-1$.\\
Dispersion & The dependence of energy magnitude on momentum disappears.\\
Fermi velocity & Positive rescaling and nonlinear deformation leave the sign unchanged.\\
Correlation length & It is not determined by a stable projection class.\\
Lieb-Robinson constants & They belong to the microscopic interaction and metric.\\
Coupling values & Many different couplings flatten to the same involution.\\
Finite-band unstable data & Stabilization allows addition of trivial bands.\\
Boundary convention & A bulk K-class alone does not choose an edge termination.\\
Interaction vertices & They lie outside the quadratic one-particle datum.\\
\bottomrule
\end{tabular}
\end{center}

\begin{proposition}[Positive rescaling is invisible]
\label{prop:rescaling}
Let $h$ be an invertible self-adjoint element of a unital C*-algebra and let
$c>0$.
Then
\[
\sgn(ch)=\sgn(h),
\qquad
p_{ch}=p_h.
\]
If $c\neq1$, the gap magnitude is multiplied by $c$.
Thus spectral flattening is not faithful to microscopic energy scales.
\end{proposition}

\begin{proof}
On $\spec(h)$,
\[
\sgn(c\lambda)=\sgn(\lambda)
\]
because $c$ is positive.
Functional calculus gives the first two identities.
The spectrum scales by $c$, so
\[
\dist(0,\spec(ch))
=
c\,\dist(0,\spec(h)).
\]
\end{proof}

This one-line counterexample is enough to rule out any claim that the
flattened K-class reconstructs the microscopic Hamiltonian.

\begin{example}[Trivial-band stabilization]
Let $h$ be gapped and let $a>0$.
Then
\[
h
\quad\text{and}\quad
h\oplus\diag(-a,a)
\]
represent the same reduced stable phase after the added occupied and empty
atomic bands are included in the reference convention.
The second system has two more bands and a new energy scale.
Stable classification is designed to forget that difference.
\end{example}

\subsection{A contract, not a reconstruction theorem}

We summarize the controlled map as
\[
(\A,h,h_{\mathrm{ref}},\mathcal S,\gamma)
\longmapsto
\kappa(h,h_{\mathrm{ref}}).
\]
Its contract is:

\begin{center}
\begin{tabular}{p{0.23\textwidth}p{0.63\textwidth}}
\toprule
Clause & Content \\
\midrule
Domain & Bounded quadratic data in a fixed C*-algebra with a Fermi gap.\\
Equivalence & Norm-continuous, symmetry-preserving gapped homotopy plus stabilization.\\
Output & A relative class in the appropriate operator K-group.\\
Guaranteed & Naturality, homotopy invariance, and stacking additivity.\\
Not guaranteed & Spectrum, dynamics, locality constants, EFT limit, or interacting completeness.\\
Inverse & No inverse is supplied.\\
\bottomrule
\end{tabular}
\end{center}

This contract matches the formal Lean representation
\begin{center}
\texttt{TopologicalPhases.Observables}.
\end{center}
That library records an Aoki comparison provider and C*-conditions without
pretending to prove the analytic theorem inside Lean's standard library.

\section{Low-energy theory is a separate map}

\subsection{Linearization near a crossing}

Suppose a Bloch Hamiltonian depends smoothly on momentum $k$ and a mass
parameter $m$.
Near an isolated crossing $(k_0,0)$, one may have
\[
h(k_0+p,m)
=
\sum_{j=1}^d v_jp_j\Gamma_j
+m\Gamma_0
+O(|p|^2+|m||p|+m^2),
\]
where the $\Gamma_j$ satisfy suitable Clifford relations.
The first-order term is a Dirac Hamiltonian.
Any scalar term proportional to the identity has already been removed by the
Fermi-level convention $\mu=0$ before the crossing subspace is expanded.

\begin{proposition}[Controlled local linearization]
\label{prop:linearization}
Let $h(k,m)$ be twice continuously differentiable in a neighborhood of
$(k_0,0)$ and suppose $h(k_0,0)=0$ after restriction to a fixed finite
crossing subspace.
Then Taylor's theorem gives
\[
h(k_0+p,m)
=
D_kh(k_0,0)[p]
+m\,\partial_mh(k_0,0)
+R(p,m),
\]
with
\[
\|R(p,m)\|
\leq C(|p|+|m|)^2
\]
on a sufficiently small neighborhood.
\end{proposition}

\begin{proof}
This is Taylor's theorem for maps from a finite-dimensional parameter space
to the finite-dimensional normed vector space of matrices.
The bound follows from a uniform bound on the second derivative on a compact
neighborhood.
\end{proof}

The proposition controls a local matrix approximation.
It does not prove convergence of an interacting lattice theory to a
relativistic quantum field theory.
Even in the quadratic setting, the crossing subspace, momentum window,
symmetry action, and neglected bands must be specified.

\subsection{Scaling limits require more data}

A genuine microscopic-to-EFT theorem would need at least:

\begin{enumerate}[label=(\alph*)]
  \item a sequence of lattice spacings or energy cutoffs;
  \item a renormalization prescription for fields and couplings;
  \item convergence of a stated class of correlation functions or operator
  algebras;
  \item control of irrelevant operators in the chosen topology;
  \item a proof that the microscopic symmetry and anomaly data descend;
  \item treatment of boundaries and defects if the target theory includes
  them.
\end{enumerate}

Condensed mathematics can organize the parameter object on which such data
vary.
It does not supply any of these estimates.

\begin{warning}[Flattening is not renormalization]
The homotopy $h\mapsto\sgn(h)$ acts on all energies in a gapped bounded
operator.
Renormalization integrates out or rescales degrees of freedom relative to an
energy or length scale.
The two maps can lead to related topological data, but they are not the same
construction.
\end{warning}

\subsection{Many-to-one behavior of effective descriptions}

Two lattice Hamiltonians can share the same linear Dirac term while differing
in quadratic dispersion, remote bands, and short-distance interactions.
For example, adding
\[
\delta h(p)=a|p|^2\Gamma_0
\]
does not change the first-order Dirac operator but changes the ultraviolet
spectrum.
Likewise, adding a high-energy atomic band changes the microscopic Hilbert
space without changing the low-energy crossing.

Thus $L$ is generally many-to-one even before applying $I$.
This is expected behavior for an effective theory.
It blocks reconstruction unless extra ultraviolet data are retained.

\subsection{From EFT data to stable classes}

Kitaev's free-fermion table \cite{Kitaev2009} and related tenfold-way
constructions \cite{Schnyder2008,Ryu2010} extract stable classes after fixing
dimension and symmetry.
Teo and Kane extend the controlled free classification to defect families
$H(k,r)$ \cite{TeoKane2010}.
These are strong results in their stated free-fermion domains.

Freed and Hopkins classify deformation classes of reflection-positive
invertible extended field theories using stable homotopy and bordism data
\cite{FreedHopkins2021}.
Their application to lattice phases assumes the existence and validity of
low-energy effective-field-theory approximations.
The bordism calculation is not, by itself, a theorem producing a local
lattice Hamiltonian.

\begin{remark}[Dimension convention]
Spatial lattice dimension is denoted $d$.
Spacetime dimension is $n=d+1$.
Bordism classifications are usually indexed by $n$, while Bloch Hamiltonians
are indexed by $d$.
The two indices must not be silently identified.
\end{remark}

\section{Disorder and the covariant observable algebra}

\subsection{A profinite disorder hull}

Let $D$ be a finite set of local labels and set
\[
\Omega=D^{\Z^d}.
\]
With the product topology, $\Omega$ is compact, Hausdorff, and profinite.
The translation action
\[
\alpha:\Z^d\curvearrowright\Omega
\]
shifts a configuration.

The covariant bulk algebra is
\[
\A_{\Omega}
=
C(\Omega)\rtimes_{\alpha}\Z^d.
\]
Because $\Z^d$ is amenable, full and reduced crossed products agree.
A magnetic field may require a twisted crossed product.
This noncommutative Brillouin-zone viewpoint underlies rigorous treatments of
disordered topological phases and the quantum Hall effect
\cite{Bellissard1994,ProdanSchulzBaldes2016}.

\subsection{Finite-range covariant Hamiltonians}

A finite-range covariant one-particle Hamiltonian can be written
schematically as
\[
h
=
\sum_{|x|\leq R} h_xu_x,
\qquad
h_x\in M_N(C(\Omega)),
\]
where the $u_x$ implement translations.
The coefficients satisfy the adjoint relations that make $h=h^*$.
The finite sum is an element of $M_N(\A_\Omega)$.

\begin{theorem}[Disordered Fermi projection]
\label{thm:disorder-projection}
Let $h=h^*\in M_N(\A_\Omega)$ and suppose
\[
0\notin\spec_{M_N(\A_\Omega)}(h).
\]
Then
\[
p=\chi_{(-\infty,0)}(h)
\]
belongs to $M_N(\A_\Omega)$ and defines
$[p]\in K_0(\A_\Omega)$.
Norm-continuous paths that remain invertible preserve this class.
\end{theorem}

\begin{proof}
Apply continuous functional calculus in the crossed-product C*-algebra.
The homotopy statement is \cref{prop:homotopy-invariance} with a one-point
parameter space.
\end{proof}

This result justifies the disorder algebra where a genuine spectral gap is
present.
It does not claim that every configuration has a gap merely because
$\Omega$ is compact.
The gap is a spectral property of the crossed-product element.

\subsection{Mobility gaps are not spectral gaps}

Disordered topological phases can remain quantized when the Fermi energy lies
in a localized regime even though it belongs to the spectrum.
That situation requires mobility-gap hypotheses, Sobolev or localization
algebras, trace estimates, and index pairings.
The discontinuous characteristic function is then handled in a more refined
regularity framework.

\begin{warning}
Theorem \ref{thm:disorder-projection} assumes a C*-spectral gap.
It must not be cited as a mobility-gap theorem.
The latter theory is deeper and has different hypotheses
\cite{Bellissard1994,ProdanSchulzBaldes2016}.
\end{warning}

\subsection{What profinite approximation does and does not prove}

Locally constant functions that factor through finite clopen quotients are
sup-norm dense in $C(\Omega)$.
This gives finite-resolution approximations to coefficient functions.
It does not imply that every K-class, spectral gap, or response coefficient
is determined by one finite quotient.

The safe condensed statement is functorial:
the profinite space $\Omega$ is a legitimate test object, covariant families
pull back along maps of disorder spaces, and their K-classes are natural when
the analytic hypotheses survive pullback.
The sheaf condition does not manufacture the crossed-product spectrum or its
gap.

\section{The SSH model as a complete controlled example}
\label{sec:ssh}

\subsection{Bloch Hamiltonian and spectrum}

The spinless SSH Bloch Hamiltonian is
\[
h(k;t_1,t_2)
=
(t_1+t_2\cos k)\sigma_x
+t_2\sin k\,\sigma_y.
\]
It anticommutes with the chiral operator $\Gamma=\sigma_z$.
Writing
\[
q(k)=t_1+t_2e^{ik},
\]
the energies are
\[
E_{\pm}(k)=\pm|q(k)|.
\]
The model descends from the dimerized chain introduced by Su, Schrieffer, and
Heeger \cite{SSH1979}.

\begin{proposition}[SSH gap]
\label{prop:ssh-gap}
For real $t_1,t_2$, the distance from zero to the spectrum is
\[
\gamma_{1/2}
=
\bigl||t_1|-|t_2|\bigr|.
\]
The separation between the occupied and empty bands is
\[
\gamma_{\mathrm{band}}
=
2\bigl||t_1|-|t_2|\bigr|.
\]
The gap closes exactly when $|t_1|=|t_2|$.
\end{proposition}

\begin{proof}
The squared positive energy is
\[
|q(k)|^2
=
t_1^2+t_2^2+2t_1t_2\cos k.
\]
Its minimum over $k$ is
\[
t_1^2+t_2^2-2|t_1t_2|
=
(|t_1|-|t_2|)^2.
\]
Taking square roots gives the half-gap.
The two bands lie symmetrically about zero, so their separation is twice that
number.
\end{proof}

\subsection{Flattening and winding}

When the model is gapped, define
\[
u(k)=\frac{q(k)}{|q(k)|}\in U(1).
\]
The flattened Hamiltonian is
\[
Q(k)
=
\begin{pmatrix}
0&\overline{u(k)}\\
u(k)&0
\end{pmatrix}.
\]

\begin{proposition}[SSH winding]
\label{prop:ssh-winding}
Assume $t_2\neq0$ and $|t_1|\neq|t_2|$.
With the convention $q(k)=t_1+t_2e^{ik}$,
\[
\nu(t_1,t_2)
=
\frac{1}{2\pi i}
\int_0^{2\pi}u(k)^{-1}\partial_ku(k)\,dk
=
\begin{cases}
1,&|t_1|<|t_2|,\\
0,&|t_1|>|t_2|.
\end{cases}
\]
\end{proposition}

\begin{proof}
The curve $q(k)$ is a circle of radius $|t_2|$ centered at the real number
$t_1$.
Multiplication by a nonzero real $t_2$ rotates the parametrized circle by
zero or $\pi$ but does not reverse its orientation.
The circle winds once counterclockwise around zero precisely when its center
lies inside its radius.
Otherwise it has winding zero.
\end{proof}

Choose a trivial reference with $|t_1|>|t_2|$.
Then the relative AIII class of a topological sample is one.
This is a concrete instance of the relative $K_1$ construction in
\cref{prop:chiral-unitary}.

\subsection{The gapless discriminant}

The global coupling discriminant is
\[
\Sigma
=
\{(t_1,t_2):|t_1|=|t_2|\}.
\]
Near the branch $t_1=t_2>0$ and momentum $k=\pi$, write
\[
m=t_1-t_2,
\qquad
k=\pi+p.
\]
Then
\[
q(\pi+p)
=
m-it_2p+O(p^2).
\]
The combined momentum-mass plane is essential.
Looking only at the coupling line misses the linking circle around the
isolated degeneracy in $(p,m)$.

\begin{proposition}[Local SSH transition charge]
\label{prop:ssh-charge}
Fix $t_2>0$ and orient the small linking circle by
\[
m=r\cos\theta,
\qquad
t_2p=r\sin\theta,
\qquad
0\leq\theta\leq2\pi.
\]
To leading order,
\[
\frac{q}{|q|}=e^{-i\theta}.
\]
Its degree is $-1$.
Reversing the circle orientation changes the sign.
The absolute charge is one and equals the magnitude of the winding jump.
\end{proposition}

\begin{proof}
Substitution gives
\[
q=m-it_2p
=
r(\cos\theta-i\sin\theta)
=
re^{-i\theta}.
\]
Normalization removes $r$.
The map $\theta\mapsto e^{-i\theta}$ has degree $-1$.
The higher-order term is uniformly smaller than $r$ on a sufficiently small
circle, so it can be removed by a homotopy that avoids zero.
\end{proof}

Topologically, the normalized map defines a class in
\[
H^1(S^1;\Z)\cong\Z.
\]
The connecting morphism for the pair $(D^2,S^1)$ sends it to
\[
\partial\nu
\in
H^2(D^2,S^1;\Z).
\]
This is the elementary model for a relative transition charge.
Teo and Kane's free-fermion defect theory gives the broader controlled
precedent \cite{TeoKane2010}.

\subsection{A symmetry-breaking detour}

Add a staggered sublattice potential
\[
\Delta\sigma_z.
\]
The Hamiltonian becomes
\[
h_{\Delta}(k)
=
(t_1+t_2\cos k)\sigma_x
+t_2\sin k\,\sigma_y
+\Delta\sigma_z.
\]
For $\Delta\neq0$, it no longer anticommutes with $\sigma_z$.

\begin{theorem}[Gapped detour after forgetting chiral symmetry]
\label{thm:ssh-detour}
Fix $t_2>0$ and $0<r<t_2$.
For $0\leq\theta\leq\pi$, set
\[
t_1(\theta)=t_2+r\cos\theta,
\qquad
\Delta(\theta)=r\sin\theta.
\]
The endpoint at $\theta=0$ is the chiral trivial phase and the endpoint at
$\theta=\pi$ is the chiral topological phase.
Along the full path, the occupied-empty band separation is exactly $2r$.
The interior breaks chiral symmetry.
\end{theorem}

\begin{proof}
Since $0<r<t_2$, both $t_1(\theta)$ and $t_2$ are positive.
The minimum positive energy is
\[
\begin{aligned}
\min_k E_+(k,\theta)
&=
\sqrt{(t_1(\theta)-t_2)^2+\Delta(\theta)^2}
\\
&=
\sqrt{r^2\cos^2\theta+r^2\sin^2\theta}
\\
&=r.
\end{aligned}
\]
Thus the full band separation is $2r$.
At the endpoints $\Delta=0$.
For $0<\theta<\pi$, $\Delta>0$ and the chiral anticommutation relation
fails.
\end{proof}

The theorem shows that the SSH charge is relative to the AIII symmetry
problem.
If the symmetry is removed from the category, the two endpoints are
gapped-homotopic.
This is not a failure of the AIII invariant.
It is a reminder that symmetry is part of its domain.

\subsection{Executable finite checks}

The accompanying Haskell program implements the exact two-band formulas.
Its checks include:

\begin{center}
\begin{tabular}{p{0.42\textwidth}p{0.44\textwidth}}
\toprule
Check & Scope \\
\midrule
SSH spectrum and analytic gap & Finite $2\times2$ Bloch matrices \\
Normalized flattened vector & Sampled momenta away from the discriminant \\
Analytic and numerical winding & AIII two-band family \\
Relative winding difference & Chosen trivial reference \\
Local linking degree & The circle $q=m-ip$ \\
Symmetry-breaking detour & A sampled semicircle with exact gap formula \\
Information-loss contract & Nonempty, pairwise-disjoint information categories \\
Three-arrow audit & Distinct typed statuses \\
\bottomrule
\end{tabular}
\end{center}

The code compiles with strict warnings and exits nonzero if a check fails.
It is not evidence for a general thermodynamic limit.

\section{Interactions: the BDI boundary}

\subsection{The free integer invariant}

Class BDI describes one-dimensional free fermions with particle-hole symmetry,
time-reversal symmetry squaring to $+1$, and their product chiral symmetry.
In the stable free classification, the phase index is an integer.
Stacking adds the integers.

This statement is different from the complex AIII SSH invariant, even though
both use winding in simple matrix representatives.
The real symmetry structure changes the classification problem.

\subsection{Fidkowski-Kitaev reduction}

Fidkowski and Kitaev constructed a symmetry-preserving interacting path that
connects free BDI phases whose integer labels differ by eight
\cite{FidkowskiKitaev2009}.
They later proved that the eight residue classes are distinct and exhaustive
for the specified one-dimensional interacting fermionic problem
\cite{FidkowskiKitaev2011}.

\begin{theorem}[BDI interaction reduction, cited]
\label{thm:bdi-reduction}
Within the one-dimensional BDI setting of
\cite{FidkowskiKitaev2009,FidkowskiKitaev2011}, the stable free index
\[
\nu_{\mathrm{free}}\in\Z
\]
descends under symmetry-preserving interactions to
\[
\nu_{\mathrm{int}}\in\Z/8.
\]
In particular, eight copies of the generating free chain admit a gapped
interacting deformation to the trivial phase.
\end{theorem}

\begin{proof}[Source-based proof sketch]
Fidkowski and Kitaev identify the boundary Majorana degrees of freedom and
construct quartic interactions that gap eight modes without breaking the BDI
symmetry.
Their explicit interacting path avoids the free critical locus while
remaining gapped.
The matrix-product-state and central-extension analysis distinguishes the
eight remaining classes and proves exhaustion.
The full proof is contained in the cited papers; it is not reproduced here.
\end{proof}

\subsection{What the reduction proves}

The reduction proves that the free classification is not a complete invariant
of the larger interacting category.
It does not invalidate the free functional-calculus theorem.
That theorem remains correct on its quadratic domain.

The categorical picture is a quotient-like comparison
\[
\Z
\longrightarrow
\Z/8
\]
after the class of allowed paths is enlarged to include interactions.
The kernel is not visible within free gapped homotopy.

\begin{corollary}[No unchanged universal extension]
There is no classification of all interacting one-dimensional BDI phases by
the free integer invariant that simultaneously agrees with free stacking and
identifies phases connected by the Fidkowski-Kitaev interacting paths.
\end{corollary}

\begin{proof}
The free invariant assigns different integers to phases whose labels differ by
eight.
The cited interacting path declares those phases equivalent.
Any invariant of interacting phase equivalence must take equal values on the
two endpoints.
Therefore the free integer cannot descend unchanged.
\end{proof}

This gives a direct counterexample to the identification of
microscopic phases with free K-theory.

\section{Realization is not the inverse map}

\subsection{Stable free representatives}

In controlled free-fermion problems, Bott periodicity and Clifford-module
constructions often provide representatives for stable K-classes
\cite{Kitaev2009,Thiang2016}.
Lattice Dirac models can realize many table entries after fixing a symmetry
class and allowing stabilization.

Even there, one must distinguish:

\begin{enumerate}[label=(\alph*)]
  \item a continuous family of finite matrices on a Brillouin torus;
  \item a finite-range or rapidly decaying real-space Hamiltonian;
  \item a chosen boundary termination;
  \item a material realization with fixed orbital constraints.
\end{enumerate}

A homotopy class of maps may have a continuous representative whose Fourier
coefficients decay but are not finite range.
Approximation can often recover finite range while preserving a positive gap,
but the approximation theorem and symmetry constraints must be stated.

\subsection{Interacting and bordism classes}

For an abstract invertible field-theory or bordism class, a microscopic
realization needs more:

\begin{enumerate}[label=(\roman*)]
  \item local degrees of freedom and an on-site or spatial symmetry action;
  \item a uniformly local or almost-local interaction;
  \item a ground-state and boundary convention;
  \item a gap uniform in system size;
  \item a computation identifying the low-energy or response class;
  \item an inverse under stacking if invertibility is claimed.
\end{enumerate}

The bordism group does not package these witnesses.
It classifies field-theoretic deformation data under its own axioms.
Freed and Hopkins explicitly frame the lattice application through an EFT
assumption \cite{FreedHopkins2021}.

\begin{definition}[Realization status]
For a candidate class $c$, assign one of the following statuses:
\begin{itemize}
  \item \status{realized}: an explicit local interaction, symmetry, uniform
  gap theorem, and comparison computation are known;
  \item \status{candidate}: a model exists, but one of those checks is
  incomplete;
  \item \status{open}: no realization or obstruction theorem is known;
  \item \status{obstructed}: a theorem rules out realization under the stated
  microscopic rules.
\end{itemize}
\end{definition}

Absence of a known model means \status{open}, not \status{obstructed}.

\subsection{Why surjectivity is hard}

Essential surjectivity of a realization map would require every abstract
class to admit a microscopic witness.
Three independent obstructions can intervene.

First, the abstract target may assume relativistic locality or full
extension to defects that a lattice construction has not supplied.
Second, anomaly constraints may force a boundary interpretation rather than
an autonomous lattice bulk in the stated dimension.
Third, a formal phase label may have no known system-size-independent gap
proof for its proposed Hamiltonian.

These are mathematical requirements, not matters of terminology.

\subsection{Why injectivity is hard}

Injectivity would say that two microscopic phases with the same abstract
class are connected by an allowed stabilized gapped path.
The target class may omit crystalline data, noninvertible excitations,
boundary conditions, or unstable finite-band information.
If any omitted datum is retained by the microscopic equivalence, injectivity
fails.

The BDI reduction shows the opposite possibility as well.
Enlarging the microscopic path category can identify free classes that were
distinct before interactions were allowed.
The answer depends on the domain and its morphisms.

\section{A condensed-mathematics formulation}

\subsection{Representable parameter objects}

For a compactly generated space $S$, write
\[
\underline S(T)=\Cont(T,S)
\]
on profinite test spaces $T$.
The precise full-faithfulness statement uses the compact-generation and
cardinal conventions of condensed mathematics
\cite{Scholze2026}.

A compact parameter family of controlled Hamiltonians is an element
\[
h\in M_N(C(S,\A))_{\mathrm{sa}}.
\]
The uniformly gapped locus consists of those $h$ for which one lower bound
works on all of $S$.
Theorem \ref{thm:flattening} gives a natural map
\[
\operatorname{Flat}_S:
\Ham^{\mathrm{free,gap}}(S)
\longrightarrow
\operatorname{Proj}(M_\infty(C(S,\A))).
\]

\begin{proposition}[Descent at the C*-algebra level]
Let $S_i\to S$ be a finite jointly surjective family of profinite spaces.
Then
\[
C(S,\A)
\longrightarrow
\operatorname{Eq}
\left(
\prod_iC(S_i,\A)
\rightrightarrows
\prod_{i,j}C(S_i\times_SS_j,\A)
\right)
\]
is an isomorphism of C*-algebras.
Compatible flattened projections glue uniquely.
\end{proposition}

\begin{proof}
Continuous functions glue uniquely across a finite jointly surjective compact
Hausdorff cover because the induced disjoint-union map is a quotient map.
Operations, involution, and the sup norm are pointwise.
The projection equation is also pointwise, so compatible local projections
glue to a projection.
\end{proof}

This is an exact statement about representable profinite tests.
It does not prove that the spectral gap itself descends from arbitrary local
gap constants unless the constants have a common positive lower bound.

\subsection{The proposed comparison diagram}

The controlled map and the broader comparison program are displayed as two
separate chains:
\[
\begin{tikzcd}[column sep=huge]
\Ham^{\mathrm{free,gap}}
  \arrow[r,"\operatorname{Flat}","\status{established}"']
&
K_{\mathrm{op}}.
\end{tikzcd}
\]
\[
\begin{tikzcd}[column sep=large]
\Ham^{\mathrm{micro,gap}}
  \arrow[r,dashed,"L","\status{open}"']
&
\EFT^{\mathrm{invertible}}
  \arrow[r,"I","\status{established}"']
&
\Class^{\mathrm{bord}}
  \arrow[r,dashed,"R","\status{open}"']
&
\Real_{\mathrm{lattice}}.
\end{tikzcd}
\]

The first arrow is the theorem proved in the controlled free sector.  In the
second chain, $I$ is solid only on the specified reflection-positive
invertible-EFT domain classified by Freed and Hopkins.  The low-energy map
$L$ and the realization map $R$ require separate comparison and realization
theorems.  No map $K_{\mathrm{op}}\to\EFT^{\mathrm{invertible}}$ and no
commutative comparison square are claimed here.

\subsection{Solidification is a K-theory bridge}

Aoki proves that solidification of algebraic K-theory recovers connective
topological K-theory for real and complex Banach algebras, with the stated
connective and Bott-periodic qualifications \cite{Aoki2024}.
This supplies a genuine condensed-to-operator K-theory bridge.

It does not reconstruct the C*-norm, positive cone, state space, or
microscopic interaction.
It also does not prove that every operator K-class is a physical phase.
The bridge is at the level of K-theory after the Banach algebra is already
given.

\section{Status ledger and formal interfaces}

\subsection{Claim ledger}

\begin{longtable}{p{0.15\textwidth}p{0.42\textwidth}p{0.31\textwidth}}
\toprule
Status & Claim & Hypotheses or source \\
\midrule
\endfirsthead
\toprule
Status & Claim & Hypotheses or source \\
\midrule
\endhead
\status{established} & Continuous functional calculus produces the flattened involution and Fermi projection. & Bounded self-adjoint C*-element with a spectral gap.\\
\status{established} & Gapped norm-continuous paths preserve the K-class. & Common gap and fixed algebra/symmetry.\\
\status{established} & The relative class is additive under block direct sum. & Fixed reference and stabilization convention.\\
\status{established} & The SSH winding is one inside the coupling circle and zero outside. & Gapped chiral two-band model.\\
\status{established} & The local SSH linking charge has magnitude one. & Fixed orientation near the isolated crossing.\\
\status{established} & A staggered mass opens a gapped detour. & Chiral symmetry is not required along the path.\\
\status{established} & BDI interactions reduce the free integer to eight classes. & Fidkowski-Kitaev one-dimensional symmetry setting.\\
\status{established} & A crossed-product Fermi projection gives a K-class. & Covariant bounded Hamiltonian with C*-spectral gap.\\
\status{proposed} & These maps assemble into a condensed Hamiltonian-stack comparison. & Representable tests are controlled; general stack construction remains.\\
\status{open} & General interacting microscopic systems are equivalent to an EFT stack. & No theorem in this scope.\\
\status{open} & Every abstract invertible bordism class has a local gapped lattice representative. & Requires explicit realization and uniform gap.\\
\status{obstructed} & Free K-theory classifies all interacting BDI phases without modification. & Counterexample from the eightfold reduction.\\
\bottomrule
\end{longtable}

\subsection{Lean representation}

The reviewed Lean foundation intentionally represents theorem interfaces, not
deep analytic proofs.
The declarations relevant here include:

\begin{itemize}
  \item \texttt{ClaimStatus} and \texttt{Provenance} for the status ledger;
  \item \texttt{FiniteVolumeGapWitness} for an explicit common lower bound;
  \item \texttt{ObjectwiseInvolution} and
  \texttt{ObjectwiseCStarNormConditions} for operator-algebraic input;
  \item \texttt{Stacking} and \texttt{PhaseEquivalence} for the stable phase
  interface;
  \item \texttt{AokiConnectiveComparisonInterface} for the imported
  solidification theorem.
\end{itemize}

No Lean declaration in that foundation asserts a Lieb-Robinson estimate,
thermodynamic limit, or Aoki comparison without a provider.
The library builds without user axioms.

\subsection{Haskell representation}

The Haskell directory
\texttt{src/microscopic-effective-theories/} contains four modules.

\begin{center}
\begin{tabular}{ll}
\toprule
Module & Role \\
\midrule
\texttt{Core.hs} & SSH spectrum, gap, flattening, winding, detour, contracts \\
\texttt{Properties.hs} & Deterministic behavioral properties \\
\texttt{Proofs.hs} & Finite equational and numerical checks \\
\texttt{Main.hs} & Demonstrations, checks, and nonzero failure exit \\
\bottomrule
\end{tabular}
\end{center}

All top-level bindings have explicit type signatures and each module has an
explicit export list.
The strict compiler flags are
\[
\texttt{-Wall -Wextra -Werror}.
\]
The program uses no claim stronger than its finite computations.

\section{Limitations and open comparison problems}

\subsection{Interacting spectral flattening}

A general interacting infinite-volume system has no bounded one-particle
Hamiltonian to which this construction can be applied.
One can flatten finite-volume matrices, but the resulting operator may be
highly nonlocal and need not converge to a quasi-local thermodynamic object.
Spectral flow is the more appropriate interacting tool, provided a uniform
gap and locality estimates have already been proved.

An open comparison problem is to relate interacting spectral-flow classes to
operator-algebraic or generalized cohomology invariants without discarding
the relevant higher defects.

\subsection{Noninvertible topological order}

A spectrum of invertible phases cannot classify general noninvertible
topological order.
Anyons, fusion, braiding, and higher defect categories require richer targets.
Two systems can share an invertible response class while having different
noninvertible excitation data.

The paper therefore restricts its bordism discussion to invertible sectors.

\subsection{Boundaries}

Bulk K-theory can pair with boundary maps and produce bulk-boundary
correspondences under controlled hypotheses
\cite{ProdanSchulzBaldes2016,TeoKane2010}.
A bulk class alone does not specify a boundary Hamiltonian, termination, or
symmetry-breaking boundary condition.
Any claimed microscopic/EFT equivalence with boundaries must include those
objects and their morphisms.

\subsection{Crystalline and spatial symmetries}

Spatial symmetries are not ordinary on-site internal symmetries.
Their action on momentum, position, disorder, and defects changes the target
K-theory or equivariant homotopy theory.
A comparison that forgets the spatial action can merge distinct crystalline
phases.

The present paper treats symmetry abstractly through a fixed C*-algebraic
structure and does not claim a complete crystalline classification.

\subsection{Mobility-gap extension}

A mobility-gap extension matters for disordered systems.
It requires localization estimates that ensure the Fermi projection belongs
to an algebra on which index pairings remain defined.
Developing that extension in condensed families would require uniform
localization control across the parameter object.
Compactness or profiniteness alone does not provide it.

\subsection{Realization tests}

A useful next step is a machine-readable realization registry.
Each row would record:

\begin{enumerate}[label=(\roman*)]
  \item the abstract class and symmetry type;
  \item an explicit microscopic interaction;
  \item a locality certificate;
  \item a uniform gap theorem;
  \item a class computation;
  \item the status \status{realized}, \status{candidate}, \status{open}, or
  \status{obstructed}.
\end{enumerate}

Such a registry would expose which part of the inverse problem is missing for
each example.

\section{Discussion}

The controlled theorem starts only after a C*-algebra has been chosen and the
Hamiltonian is bounded and
self-adjoint, the Fermi level is uniformly isolated, and the symmetry and
reference are fixed.
Under those assumptions, functional calculus gives a stable class with clean
naturality properties.

The loss is built into stable K-theory, which forgets deformations and added
trivial bands.
Describing that quotient as a reconstruction of the microscopic system would
be incorrect.

The SSH calculation makes both sides visible.
Within AIII symmetry, the winding jump forces a gap closing and the linking
circle has charge one in magnitude.
Once the staggered mass is permitted, the detour keeps a constant gap.
The category changed, so the classification changed.

The BDI example makes the same point at the interaction boundary.
Allowing quartic interactions enlarges the path space and identifies every
eighth free phase.
No amount of categorical packaging removes that physical fact.

Condensed mathematics enters most cleanly through parameter functoriality and
the treatment of profinite disorder spaces.
It places continuous families, inverse limits, and descent in one formal
setting.
Its contribution is organizational and homological.
The analytic inputs remain locality, positivity, spectral control, and the
existence of comparison maps.

\section{Conclusion}

There is a rigorous bridge from a controlled free-fermion Hamiltonian to a
relative operator K-class:
\[
h
\longmapsto
\sgn(h)
\longmapsto
p_h
\longmapsto
[p_h]-[p_{h_{\mathrm{ref}}}].
\]
The bridge is natural in compact parameter spaces, invariant under uniformly
gapped norm-continuous homotopy, and additive under stacking.
For chiral systems it has the corresponding odd-unitary form.
For spectrally gapped covariant disorder models it lives in the K-theory of a
crossed product.

This bridge is not an equivalence between microscopic lattice systems and
effective field theories.
It forgets energy scales and much of the ultraviolet model.
It does not include general interactions.
It has no automatic inverse realization map.

The correct research program therefore keeps the three arrows separate.
One proves a low-energy limit where possible, extracts an invariant under
fixed hypotheses, and treats microscopic realization as its own theorem or
open problem.
The SSH transition and the BDI interaction reduction show why those
distinctions are physically necessary.

\appendix

\section{Functional-calculus details}

\subsection{Continuity of the sign map on a gapped set}

Fix $\gamma>0$ and $M>\gamma$.
On the compact set
\[
X_{\gamma,M}
=
[-M,-\gamma]\cup[\gamma,M],
\]
the sign function is continuous.
By the Stone-Weierstrass theorem, polynomials can approximate it uniformly on
$X_{\gamma,M}$.
If $h$ is self-adjoint with
$\spec(h)\subset X_{\gamma,M}$, then
$\sgn(h)$ is the norm limit of the same polynomials in $h$.

For a norm-continuous family with a common gap and a common norm bound, the
polynomial approximants are uniformly continuous in the family parameter.
Their uniform limit is norm-continuous.
This gives an elementary route to the continuity assertion in
\cref{prop:homotopy-invariance}.

\subsection{A resolvent formula}

The negative spectral projection can also be written as a Riesz projection.
Choose a contour $\mathcal C$ enclosing the negative spectrum and excluding
the positive spectrum.
Then
\[
p_h
=
\frac{1}{2\pi i}
\int_{\mathcal C}(z-h)^{-1}\,dz.
\]
If $h'$ is close enough to $h$, the same contour lies in the resolvent set of
$h'$.
The resolvent identity
\[
(z-h')^{-1}-(z-h)^{-1}
=
(z-h')^{-1}(h'-h)(z-h)^{-1}
\]
gives norm continuity of the projection.

This formula also displays the role of the gap.
As the contour approaches the spectrum, the resolvent norm grows and uniform
control is lost.

\subsection{Symmetry functoriality}

Suppose an automorphism or anti-automorphism $\alpha$ encodes a symmetry and
$h$ obeys a relation such as
\[
\alpha(h)=h,
\qquad
\alpha(h)=-h,
\qquad\text{or}\qquad
\alpha(h(k))=h(-k).
\]
Functional calculus transports the corresponding relation to $q=\sgn(h)$
because the sign function is real and odd.
The target K-group must still remember whether $\alpha$ is linear,
antilinear, grading-preserving, or grading-reversing.

\section{Difference classes and references}

\subsection{Why a reference is useful}

The Grothendieck group of projections contains formal differences.
Physically, choosing $h_{\mathrm{ref}}$ identifies which atomic or vacuum
configuration counts as zero.
The class
\[
[p_h]-[p_{h_{\mathrm{ref}}}]
\]
then measures an obstruction between two gapped data rather than assigning an
absolute label to one Hamiltonian.

Thiang's formulation makes this relative viewpoint explicit
\cite{Thiang2016}.
It avoids suggesting that the abstract group element contains a preferred
microscopic representative.

\subsection{Cocycle law}

For three controlled systems $h_0,h_1,h_2$ in the same target,
\[
\kappa(h_0,h_2)
=
\kappa(h_0,h_1)+\kappa(h_1,h_2),
\]
where
\[
\kappa(h_i,h_j)=[p_{h_i}]-[p_{h_j}].
\]
The formula is immediate by cancellation in the abelian group.
It is checked on finite SSH representatives by the Haskell program.

\subsection{Reference changes}

If the reference is changed from $r$ to $r'$, then
\[
\kappa(h,r')
=
\kappa(h,r)+\kappa(r,r').
\]
Thus absolute coordinates change by a constant translation.
Differences between two physical systems do not depend on that coordinate
choice.

\section{SSH conventions and numerical checks}

\subsection{Orientation}

With
\[
q(k)=t_1+t_2e^{ik},
\]
the topological SSH winding is $+1$ for $|t_1|<|t_2|$.
Near $k=\pi$, the local coordinate is
\[
q\simeq m-it_2p.
\]
The positively oriented circle $(m,t_2p)=(r\cos\theta,r\sin\theta)$ maps to
$e^{-i\theta}$ and therefore has degree $-1$.
There is no contradiction.
The local charge and the phase jump use different boundary orientations.
Only the orientation-adjusted equality, or equality in magnitude, is
canonical.

\subsection{Discrete winding computation}

For nonzero samples $q_j=q(k_j)$ with
$k_j=2\pi j/N$, the code computes
\[
\nu_N
=
\operatorname{round}
\left(
\frac{1}{2\pi}
\sum_{j=0}^{N-1}
\operatorname{Arg}
\left(\frac{q_{j+1}}{q_j}\right)
\right).
\]
The analytic gap check runs first, so division by zero is excluded.
Set $g_\chi=||t_1|-|t_2||$ and $h=2\pi/N$.
Since
\[
\left|\frac{d}{dk}\operatorname{Arg}q(k)\right|
\leq \frac{|t_2|}{|q(k)|}
\leq \frac{|t_2|}{g_\chi},
\]
the code requires $h|t_2|/g_\chi<\pi$ before summing principal phase
increments.
This sufficient sampling condition prevents a true increment from crossing
the branch cut; an explicit runtime check rejects an increment of magnitude
$\pi$ as well.
The result then agrees with \cref{prop:ssh-winding}.

The computation is a regression test for the implementation.
It is not used as the proof of the analytic proposition.

\subsection{Finite gap sampling}

The executable proof check compares the closed formula
\[
2\sqrt{(|t_1|-|t_2|)^2+\Delta^2}
\]
with the minimum sampled band separation on a mesh that contains $k=\pi$.
The analytic derivation remains the proof.
The mesh catches sign, factor-of-two, and convention mistakes in code.

\section{Comparison checklist for later work}

Before a future paper claims a microscopic/EFT equivalence, it should fill in
the following record.

\begin{longtable}{p{0.25\textwidth}p{0.61\textwidth}}
\toprule
Field & Required answer \\
\midrule
\endfirsthead
\toprule
Field & Required answer \\
\midrule
\endhead
Microscopic objects & Interactions, Hilbert spaces, symmetries, and boundary conventions.\\
Microscopic morphisms & Exact definition of gapped path, circuit, or quasi-local equivalence.\\
Effective objects & Fields, tangential structure, anomalies, and defect depth.\\
Effective morphisms & Deformation, natural equivalence, or bordism convention.\\
Limit & Scaling parameter and topology of convergence.\\
Uniformity & Bounds uniform in volume and external parameters.\\
Construction & Formula or theorem defining the forward map.\\
Invariance & Proof that microscopic equivalences map to effective equivalences.\\
Faithfulness & Proof or counterexample.\\
Fullness & Proof that effective morphisms lift.\\
Surjectivity & Realization theorem for every target class.\\
Stacking & Compatibility with direct sum or tensor product.\\
Symmetry & Preservation of internal, antiunitary, and spatial actions.\\
Boundary & Treatment of edges, defects, and anomaly inflow.\\
Lost data & Explicit list of ultraviolet information removed by the map.\\
Interaction boundary & Statement of which nonquadratic perturbations are admitted.\\
\bottomrule
\end{longtable}

Most known general comparisons fill only part of this record.
That is still mathematically useful.
The problem arises when a partial invariant-producing map is reported as a
categorical equivalence.

\begin{thebibliography}{99}

\bibitem{Aoki2024}
K.~Aoki,
\emph{(Semi)topological K-theory via solidification},
arXiv:2409.01462, 2024.

\bibitem{Atiyah1967}
M.~F.~Atiyah,
\emph{K-Theory},
W.~A.~Benjamin, 1967.

\bibitem{Bellissard1994}
J.~Bellissard, A.~van Elst, and H.~Schulz-Baldes,
\emph{The non-commutative geometry of the quantum Hall effect},
Journal of Mathematical Physics \textbf{35} (1994), pages 5373 to 5451,
\href{https://doi.org/10.1063/1.530758}{doi:10.1063/1.530758}.

\bibitem{Blackadar1998}
B.~Blackadar,
\emph{K-Theory for Operator Algebras}, second edition,
Cambridge University Press, 1998.

\bibitem{Scholze2026}
P.~Scholze,
\emph{Lectures on Condensed Mathematics},
arXiv:2605.03658, 2026.

\bibitem{FidkowskiKitaev2009}
L.~Fidkowski and A.~Kitaev,
\emph{The effects of interactions on the topological classification of free
fermion systems},
Physical Review B \textbf{81} (2010), 134509,
\href{https://doi.org/10.1103/PhysRevB.81.134509}{doi:10.1103/PhysRevB.81.134509}.

\bibitem{FidkowskiKitaev2011}
L.~Fidkowski and A.~Kitaev,
\emph{Topological phases of fermions in one dimension},
Physical Review B \textbf{83} (2011), 075103,
\href{https://doi.org/10.1103/PhysRevB.83.075103}{doi:10.1103/PhysRevB.83.075103}.

\bibitem{FreedHopkins2021}
D.~S.~Freed and M.~J.~Hopkins,
\emph{Reflection positivity and invertible topological phases},
Geometry \& Topology \textbf{25} (2021), pages 1165 to 1330,
\href{https://doi.org/10.2140/gt.2021.25.1165}{doi:10.2140/gt.2021.25.1165}.

\bibitem{Kitaev2001}
A.~Kitaev,
\emph{Unpaired Majorana fermions in quantum wires},
Physics-Uspekhi \textbf{44} (2001), pages 131 to 136,
\href{https://doi.org/10.1070/1063-7869/44/10S/S29}{doi:10.1070/1063-7869/44/10S/S29}.

\bibitem{Kitaev2009}
A.~Kitaev,
\emph{Periodic table for topological insulators and superconductors},
AIP Conference Proceedings \textbf{1134} (2009), pages 22 to 30,
\href{https://doi.org/10.1063/1.3149495}{doi:10.1063/1.3149495}.

\bibitem{ProdanSchulzBaldes2016}
E.~Prodan and H.~Schulz-Baldes,
\emph{Bulk and Boundary Invariants for Complex Topological Insulators},
Springer, 2016,
\href{https://doi.org/10.1007/978-3-319-29351-6}{doi:10.1007/978-3-319-29351-6}.

\bibitem{Ryu2010}
S.~Ryu, A.~P.~Schnyder, A.~Furusaki, and A.~W.~W.~Ludwig,
\emph{Topological insulators and superconductors: tenfold way and
dimensional hierarchy},
New Journal of Physics \textbf{12} (2010), 065010,
\href{https://doi.org/10.1088/1367-2630/12/6/065010}{doi:10.1088/1367-2630/12/6/065010}.

\bibitem{Schnyder2008}
A.~P.~Schnyder, S.~Ryu, A.~Furusaki, and A.~W.~W.~Ludwig,
\emph{Classification of topological insulators and superconductors in three
spatial dimensions},
Physical Review B \textbf{78} (2008), 195125,
\href{https://doi.org/10.1103/PhysRevB.78.195125}{doi:10.1103/PhysRevB.78.195125}.

\bibitem{SSH1979}
W.~P.~Su, J.~R.~Schrieffer, and A.~J.~Heeger,
\emph{Solitons in polyacetylene},
Physical Review Letters \textbf{42} (1979), pages 1698 to 1701,
\href{https://doi.org/10.1103/PhysRevLett.42.1698}{doi:10.1103/PhysRevLett.42.1698}.

\bibitem{TeoKane2010}
J.~C.~Y.~Teo and C.~L.~Kane,
\emph{Topological defects and gapless modes in insulators and
superconductors},
Physical Review B \textbf{82} (2010), 115120,
\href{https://doi.org/10.1103/PhysRevB.82.115120}{doi:10.1103/PhysRevB.82.115120}.

\bibitem{Thiang2016}
G.~C.~Thiang,
\emph{On the K-theoretic classification of topological phases of matter},
Annales Henri Poincare \textbf{17} (2016), pages 757 to 794,
\href{https://doi.org/10.1007/s00023-015-0418-9}{doi:10.1007/s00023-015-0418-9}.

\end{thebibliography}

\end{document}
